% Full proof of the nine-operation replacement for the signed-index block.
\section{A smaller auxiliary Pell parameter: 92 operations}
\label{sec:half-parameter-pell}

The square $K=(ic^2)^2$ already present in the relaxed auxiliary norm
allows a smaller Pell parameter. Replacing the ten-operation signed
block by a nine-operation block saves one multiplication. The new
block uses one additional positive unknown and one additional equation.
Its positive witnesses are constructed afresh; the coefficient encoding
and fixed index do not change.

\begin{theorem}\label{thm:best92}
Every recursively enumerable subset of the positive integers has an
admissible fixed index $(V,H,T_L)$ for a system with 35 positive
unknowns and 23 equations admitting a straight-line certificate of
92 operations: 49 multiplications and 43 additions. The fixed index
and coefficient encoding are those of Theorem~\ref{thm:best93}.
The supplied inputs are $(x,V,H,T_L)$, with three fixed components
besides the positive input $x$. Fixed numerals and equality tests
are free.
\end{theorem}

\subsection{The new equations}

Retain the definitions
\begin{equation}\label{eq:half-pell-parameters}
 J=2r+1,\qquad A=a+4,\qquad D=A^2-1,\qquad
 R=ic^2,\qquad K=R^2.
\end{equation}
In particular, the relaxed auxiliary equation is retained:
\begin{equation}\label{eq:half-pell-relaxed}
 R^2=D(f^2-1).
\end{equation}
The former signed equation was
\[
 v(v+1)=K(K-1)(J+jc)^2,\qquad v=of-d^2.
\]
Replace it by
\begin{equation}\label{eq:half-pell-linear}
 u=J+jc,\qquad u=c+of,
\end{equation}
and
\begin{equation}\label{eq:half-pell-norm}
 K(u^2-y_{\rm aux}^2)=1-y_{\rm aux}^2.
\end{equation}
The quantity $u$ is an arithmetic register, not another supplied
unknown. The old variables $o,j$ and the new variable $y_{\rm aux}$
are positive integers. The last equation is equivalently
\begin{equation}\label{eq:half-pell-standard-norm}
 (Ru)^2-(R^2-1)y_{\rm aux}^2=1.
\end{equation}
Its ordinary Pell parameter is $R$, whereas the preceding signed
construction used $2R^2-1$. The certificate does not have to
construct either $Ru$ or that larger parameter.

\subsection{Independent bounds before the new block}

The preliminary proof of Section~\ref{sec:affine-radix}, retained
by the subsequent arithmetic improvements, gives
\begin{equation}\label{eq:half-pell-preliminary}
 \begin{gathered}
 n=q^8,\qquad n\ge64,\qquad n\le r<2n^3,\qquad 3L\le B\le n,\\
 U=wn^2,\qquad Y_P=s_Pn^2,\qquad a=Y_P(U+1),\qquad
 P=2UY_P^2+1,\\
 U,Y_P\ge n^2,\qquad UY_P>r+1,\qquad a>J,\qquad c>J>1.
 \end{gathered}
\end{equation}
The subscripts distinguish these Pell quantities from the encoded
logical coordinates. These bounds use positivity, the combined gap,
geometry, packing, and the positive interval. They precede coefficient
decoding and any use of either auxiliary signed norm.

Classification of the unchanged main and first norms gives positive
indices $p,t$ with
\[
 c=\psi_A(p),\qquad d=\chi_A(p),\qquad k=\psi_P(t).
\]
The first-index equation implies $t\ge r+1$. Since $P>A$ and
$c>Y_Pk>k$, the growth comparison in
Section~\ref{sec:relaxed-auxiliary} forces
\begin{equation}\label{eq:half-pell-large-index}
 p\ge r+2\ge66,\qquad c>AD^2.
\end{equation}
The retained relaxed-norm argument therefore applies to
\eqref{eq:half-pell-relaxed}. Its integer-coefficient fundamental-unit
and divisibility proof supplies a positive index $m$ such that
\begin{equation}\label{eq:half-pell-auxiliary-index}
 f=\chi_A(m),\qquad p\mid m,\qquad c\mid m,
 \qquad R=D\psi_A(m).
\end{equation}
That proof uses neither the former nor the new signed-index equation.
Elementary Pell growth now gives
\begin{equation}\label{eq:half-pell-comparison}
 0<2p\le c\le m,\qquad R=ic^2\ge c^2>A>1.
\end{equation}
In particular, $R$ is an integer Pell parameter greater than one.
No exact value of $p$, no parity of $r$, no conclusion $q=B^L$,
and no decoded circuit equation has been assumed here.

\subsection{The odd-index polynomial identity}

For the standard Pell sequences, write
\[
 \begin{gathered}
 \chi_X(0)=1,\quad\chi_X(1)=X,\qquad
 \psi_X(0)=0,\quad\psi_X(1)=1,\\
 z(n+2)=2Xz(n+1)-z(n).
 \end{gathered}
\]
For every $h\ge0$ there is an integer polynomial $Q_h(T)$ with
\begin{equation}\label{eq:half-pell-Q-definition}
 \chi_X(2h+1)=XQ_h(X^2).
\end{equation}
Its initial values and doubled-step recurrence are
\begin{equation}\label{eq:half-pell-Q-recurrence}
 \begin{gathered}
 Q_0(T)=1,\qquad Q_1(T)=4T-3,\\
 Q_{h+2}(T)=(4T-2)Q_{h+1}(T)-Q_h(T).
 \end{gathered}
\end{equation}
The doubled recurrence follows by applying the ordinary recurrence
twice, or from $\chi_X(2)=2X^2-1$ and the addition identities.
Thus these formulas prove divisibility by $X$ as an integer-polynomial
identity; they do not require division modulo an integer.

The odd subsequence $\psi_A(2h+1)$ has recurrence coefficient
$4A^2-2$ and initial values $1,4A^2-1$. Consequently
$(-1)^h\psi_A(2h+1)$ has coefficient $2-4A^2$ and initial
values $1,1-4A^2$. Substituting $T=1-A^2$ in
\eqref{eq:half-pell-Q-recurrence} gives exactly these data.
Induction proves
\begin{equation}\label{eq:half-pell-Q-complement}
 Q_h(1-A^2)=(-1)^h\psi_A(2h+1).
\end{equation}
Similarly, substitution of $T=0$ gives
\begin{equation}\label{eq:half-pell-Q-zero}
 Q_h(0)=(-1)^h(2h+1).
\end{equation}
Both identities hold for every nonnegative $h$. Their two uses below
have different moduli.

\subsection{Sufficiency: the exact main index}

Suppose the new source equations have positive integer witnesses.
By \eqref{eq:half-pell-standard-norm}, positivity, and $R>1$,
ordinary Pell classification gives a positive index $s$ such that
\begin{equation}\label{eq:half-pell-classification}
 Ru=\chi_R(s),\qquad y_{\rm aux}=\psi_R(s).
\end{equation}
The index $s$ is odd. Indeed, the chi recurrence at $R=0$ gives
$\chi_R(2h)\equiv(-1)^h\pmod R$. An even $s$ would make
$R>1$ divide one, since $R$ divides the first coordinate in
\eqref{eq:half-pell-classification}. Write $s=2h+1$.
Equation~\eqref{eq:half-pell-Q-definition} gives the exact equality
\begin{equation}\label{eq:half-pell-u-polynomial}
 u=Q_h(R^2).
\end{equation}

Reduction of \eqref{eq:half-pell-relaxed} modulo $f$ gives
\[
 R^2=D(f^2-1)\equiv1-A^2\pmod f.
\]
Equations \eqref{eq:half-pell-Q-complement},
\eqref{eq:half-pell-u-polynomial}, and the second equality in
\eqref{eq:half-pell-linear} therefore imply
\[
 (-1)^h\psi_A(s)\equiv u\equiv c=\psi_A(p)\pmod f.
\]
Squaring and using the doubling identity
\[
 \chi_A(2z)=1+2D\psi_A(z)^2
\]
gives
\begin{equation}\label{eq:half-pell-chi-comparison}
 \chi_A(2s)\equiv\chi_A(2p)\pmod{\chi_A(m)}.
\end{equation}

Apply Lemma~\ref{lem:signed-chi-new}: for $A>1$ and
$0<k\le m$, the congruence
$\chi_A(n)\equiv\pm\chi_A(k)\pmod{\chi_A(m)}$ implies
$n\equiv\pm k\pmod{2m}$. The positive-sign case also follows
from the ordinary chi step-down statement with its stronger
modulus $4m$. Here its comparison-index hypothesis is exactly
$0<2p\le m$ from \eqref{eq:half-pell-comparison}. Hence
\begin{equation}\label{eq:half-pell-stepdown}
 2s\equiv\pm2p\pmod{2m},\qquad
 s\equiv\pm p\pmod m,\qquad s\equiv\pm p\pmod c.
\end{equation}
The middle implication divides an integer equality
$2s=\pm2p+2mz$ by two. It does not assume that two is invertible
in an even residue ring.

Independently, $c\mid R$ because $R=ic^2$. Equations
\eqref{eq:half-pell-Q-zero} and
\eqref{eq:half-pell-u-polynomial} give
\[
 u\equiv(-1)^h s\pmod c.
\]
The first equality in \eqref{eq:half-pell-linear}, together with
\eqref{eq:half-pell-stepdown}, now yields
\begin{equation}\label{eq:half-pell-main-congruence}
 J\equiv\pm p\pmod c.
\end{equation}
Both $J$ and $p$ lie strictly between zero and $c$, by
\eqref{eq:half-pell-preliminary} and Pell growth. Thus either
$J=p$ or $J+p=c$. The latter is impossible: the psi recurrence
gives $\psi_A(p)\equiv p\pmod2$ for every integer $A$, whereas
$J=2r+1$ is odd. Therefore $J+p$ has parity opposite to
$c=\psi_A(p)$. We have proved
\begin{equation}\label{eq:half-pell-main-index}
 p=J,\qquad c=\psi_A(J),\qquad d=\chi_A(J).
\end{equation}
This deduction did not assume that $r$ is even. In particular, it
does not use the subsequent binary decoding circularly.

\subsection{The remaining sufficiency proof}

Once \eqref{eq:half-pell-main-index} is available, the retained
first-index and ratio arguments recover the first index $r+1$,
prove $Y_P\ge U^r$ and $a>U^{r+1}$, and use the common-witness
exponent equation to obtain $U=4^J$. The unchanged fixed value
$T_L=\psi_4(L)$, gap, and second Pell equations recover the second
index $L$ and $q=B^L$. The same bounds give
$n^2\mid\binom{2r}{r}$.

These are exactly the hypotheses and conclusions of
Section~\ref{sec:affine-radix}. The replaced auxiliary variables
$o,j$ occur nowhere in these later equations. Since $n^2\mid U$,
the integers $q$ and then $B$ are powers of two. Neither $b$ nor
$H$ is required to be a power of two. The unchanged masks recover
the canonical codes, the three-way physical coordinates, all paired
rows, and the unit $\delta=1$. They therefore recover the represented
system at its actual positive input $x$.

The supplied fixed index remains
\[
 H=H_0-3,\qquad V=\ell_0(4)+e_0(4)4^L,\qquad
 T_L=\psi_4(L),\qquad B=H+b+4.
\]
The new block adds no index parameter, compiler convention, or
support requirement.

\subsection{Necessity: parity and positive auxiliary witnesses}

Start with a solution of the represented Diophantine system. Choose
its canonical physical encoding and all positive coding and Pell
witnesses as in Section~\ref{sec:affine-radix}, translating the
fixed index and retaining the factored mask as in
Sections~\ref{sec:shifted-affine} and \ref{sec:factored-mask}.
All those choices remain available.

The canonical indicator has zero unit digit, so $B\mid\ell$.
The integer $g$ also has zero unit digit, so $B\mid g$.
Since $B$ and $q=B^L$ are even, both packed integers
\[
 \begin{aligned}
 S&=g+q^2(\ell+eq+q^2\sigma),\\
 T+1&=q^2(1+\theta\lambda)+\ell(\theta q^4-b)
 \end{aligned}
\]
are even. Hence $r=S(n^2-n)+(T+1)(n^2-1)$ is even, and
\begin{equation}\label{eq:half-pell-necessity-parity}
 J=2r+1\equiv1\pmod4.
\end{equation}

Keep the canonical relaxed auxiliary witnesses
\begin{equation}\label{eq:half-pell-necessity-auxiliary}
 \begin{gathered}
 m=2cJ,\qquad f=\chi_A(m),\qquad
 i_{\rm old}=\frac{\psi_A(m)}{c^2}>0,\\
 i=Di_{\rm old},\qquad R=ic^2=D\psi_A(m),\qquad K=R^2.
 \end{gathered}
\end{equation}
The quotient is integral: expanding
$(d+c\sqrt D)^{2c}$, its term with one square root contains
$2c^2$, while every later square-root term contains at least $c^3$.
Positivity follows from $m>0$. The Pell norm gives
\eqref{eq:half-pell-relaxed}. Thus the same $i,f$ can be retained.
Keep every other canonical witness except $o,j$, and define
\begin{equation}\label{eq:half-pell-new-witnesses}
 \begin{gathered}
 y_{\rm aux}=\psi_R(J),\qquad
 u=\frac{\chi_R(J)}R=Q_{(J-1)/2}(R^2),\\
 o=\frac{u-c}{f},\qquad j=\frac{u-J}{c}.
 \end{gathered}
\end{equation}
The division defining $u$ is exact by
\eqref{eq:half-pell-Q-definition}. Equation
\eqref{eq:half-pell-necessity-parity} makes $(J-1)/2$ even, so
\eqref{eq:half-pell-Q-complement} and
\eqref{eq:half-pell-Q-zero} give
\[
 u\equiv\psi_A(J)=c\pmod f,\qquad u\equiv J\pmod c.
\]
Both quotients in \eqref{eq:half-pell-new-witnesses} are therefore
integers.

They are strictly positive. For $R>1$, the positive chi sequence
is strictly increasing, and its recurrence gives
\[
 \chi_R(k+1)>(2R-1)\chi_R(k)\qquad(k\ge1).
\]
Since $\chi_R(1)=R$ and $J\ge3$, we obtain
\begin{equation}\label{eq:half-pell-positive-quotients}
 \frac{\chi_R(J)}R>(2R-1)^{J-1}
   >(2A)^{J-1}>\psi_A(J)=c>J.
\end{equation}
The middle inequality uses the integer bound $R\ge c^2>A$,
which implies $2R-1>2A$; the next is the usual strict psi upper
bound at index $J\ge3$. Thus $o,j>0$, and $y_{\rm aux}>0$ too.

The Pell equation for $(\chi_R(J),\psi_R(J))$ is precisely
\eqref{eq:half-pell-standard-norm}, hence
\eqref{eq:half-pell-norm}. The definitions give both equalities
in \eqref{eq:half-pell-linear}. Every other source equation is
retained and uses none of the changed auxiliary values. All 35
supplied unknowns are therefore positive integers. This proves
necessity without asserting an identity map for the auxiliary
coordinates, and shows exactly why $J\equiv1\pmod4$ was needed
only in this direction.

\subsection{The complete arithmetic count}

The former signed block used ten operations:
\[
 \begin{aligned}
 O&=of,& v&=O-d^2,& v_1&=v+1,& L&=vv_1,\\
 K_1&=K-1,& K_2&=KK_1,\\
 J_c&=jc,& u&=J+J_c,& U_2&=u^2,& P&=K_2U_2.
 \end{aligned}
\]
The quantities $d^2,K,c,J$ already exist outside this block.
Replace it by nine operations:
\begin{equation}\label{eq:half-pell-nine-operations}
 \begin{aligned}
 J_c&=jc,& u&=J+J_c,& O&=of,\\
 u_{\rm rhs}&=c+O,& U_2&=u^2,& Y_2&=y_{\rm aux}^2,\\
 G&=U_2-Y_2,& L&=KG,& P&=1-Y_2.
 \end{aligned}
\end{equation}
The free tests are $u=u_{\rm rhs}$ and $L=P$. The old block
has six multiplications and four additions or subtractions; the new
block has five multiplications and four additions or subtractions.
Every other primitive remains unchanged. The complete count is
\[
 93-1=92=49\text{ multiplications}+43\text{ additions}.
\]
The variable $d$ and its squared register remain in the main norm;
their absence from the new auxiliary block gives no further saving.

The new intermediates are ordinary integer registers. In fact $G$
and $P$ in \eqref{eq:half-pell-nine-operations} are negative on
solutions. Since $u=J+jc>1$, Pell index one is excluded in
\eqref{eq:half-pell-classification}; thus $y_{\rm aux}>1$ and
\[
 u^2-y_{\rm aux}^2=\frac{1-y_{\rm aux}^2}{K}<0.
\]
These subtractions are represented by reversed addition equations.
No positivity condition is imposed on the intermediate gaps;
all 35 supplied unknowns, including $y_{\rm aux}$, are positive.

The checker writes the source norm with $K_{\rm source}=D(f^2-1)$
and uses the already computed $K=(ic^2)^2$ in the primitive product.
The difference of its residuals is exactly
\[
 (K-K_{\rm source})(u^2-y_{\rm aux}^2).
\]
Its first factor is the independently exact retained equation
\eqref{eq:half-pell-relaxed}, so this is an acyclic correction.
It replaces the preceding signed norm's quadratic-factor correction.
The new linear congruence check is exact, and the geometric and
second-exponent corrections remain as before.

Appendix~\ref{app:92} gives the complete table of 92 instructions
and all 23 equality tests. The checker is in
\file{../verification/}:
\begin{center}
\small\file{round35\_1980\_half\_parameter\_pell\_certificate.py}.
\end{center}
Its JSON receipt records the full primitive schedule, domains, source
residuals, and remaining corrections. Those arithmetic checks and the
positive-witness proof above establish Theorem~\ref{thm:best92}.
The number 92 is an upper bound; no optimality assertion or new
compiled Lean theorem is claimed.
